% \newcommand{\zdh}[1]{\textcolor{brown}{ZDH: #1}}
\newcommand{\yz}[1]{\textcolor{blue}{YZ: #1}}
% \newcommand{\rebuttal}[1]{{\color{blue}{#1}}}
% \newcommand{\rebuttal}[1]{{#1}}

\newcommand{\blue}[1]{\textcolor{blue}{#1}}
\newcommand{\med}[1]{\textcolor{magenta}{#1}}
\newcommand{\red}[1]{\textcolor{red}{#1}}
\newcommand{\gray}[1]{\textcolor{gray}{#1}}
\newcommand{\green}[1]{\textcolor{green}{#1}}

\usepackage{xspace}
\newcommand*{\eg}{\textit{e.g.}\xspace}
\newcommand*{\ie}{\textit{i.e.}\xspace}
\newcommand*{\cf}{\textit{cf.}\xspace}
%%%%%%%%%%%%%%%%%%%%%%%%%%%%%%%%
% THEOREMS
%%%%%%%%%%%%%%%%%%%%%%%%%%%%%%%%
\theoremstyle{plain}
\newtheorem{theorem}{Theorem}%[section]
\newtheorem{proposition}[theorem]{Proposition}
\newtheorem*{proposition*}{Proposition}
\newtheorem{lemma}[theorem]{Lemma}
\newtheorem{corollary}[theorem]{Corollary}
\theoremstyle{definition}
\newtheorem{definition}[theorem]{Definition}
\newtheorem{assumption}[theorem]{Assumption}
% \theoremstyle{remark}
\theoremstyle{definition}
\newtheorem{remark}[theorem]{Remark}


% modifed \bar{}
\newcommand{\overbar}[1]{\mkern 1.5mu\overline{\mkern-1.5mu#1\mkern-1.5mu}\mkern 1.5mu}

% \DeclareMathOperator*{\argmax}{arg\,max}
% \DeclareMathOperator*{\argmin}{arg\,min}

% \newcommand{\refe}{\text{ref}} % \ref is used by latex default

\newcommand*{\dif}{\mathop{}\!\mathrm{d}}
% \newcommand{\KL}[2]{\mathcal{D}_{\mathrm{KL}}\left(#1\Vert #2\right)}
\newcommand{\tr}[1]{\text{tr}\left(#1\right)}
\newcommand{\norm}[1]{\left\lVert#1\right\rVert}
\def\viz{\emph{viz}\onedot}
\newcommand{\diag}[1]{\text{diag}\left(#1\right)}
\newcommand{\dt}{\frac{\dif}{\dif t}}

\newcommand{\aaa}{\mathbf{a}}
% \newcommand{\ccc}{\mathbf{c} % weird
\newcommand{\x}{\mathbf{x}}
\newcommand{\y}{\mathbf{y}}
\newcommand{\z}{\mathbf{z}}
\newcommand{\s}{\mathbf{s}}
\newcommand{\g}{\mathbf{g}}
\newcommand{\f}{\mathbf{f}}
\newcommand{\w}{\mathbf{w}}
\newcommand{\0}{\mathbf{0}}

\newcommand{\dd}{{\boldsymbol{\delta}}}
\newcommand{\eps}{{\boldsymbol{\epsilon}}}
\newcommand{\vtheta}{{\boldsymbol{\theta}}}
\newcommand{\vphi}{{\boldsymbol{\phi}}}
\newcommand{\ra}{{\rightarrow}}
\newcommand{\disteq}{{\stackrel{d}{=}}}

% \newcommand{\W}{\mathbf{W}}
% \newcommand{\Q}{\mathcal{Q}}
\newcommand{\A}{\mathcal{A}}
% \newcommand{\gB}{\mathcal{B}}
% \newcommand{\D}{\mathfrak{D}}
\newcommand{\N}{\mathcal{N}}
\newcommand{\E}{\mathbb{E}}
\newcommand{\gE}{\mathcal{E}}
\newcommand{\LL}{{\mathcal{L}}}
\newcommand{\LDB}{\LL_{\text{DB}}}
\newcommand{\lDB}{\ell_{\text{DB}}}
\newcommand{\lFL}{\ell_{\text{FL}}}
\newcommand{\lFLDB}{\ell_{\text{FL-DB}}}
\newcommand{\lST}{\ell_{\text{SubTB}}}
\newcommand{\PP}{\mathbb{P}}
\newcommand{\R}{\mathbb{R}}
\newcommand{\I}{\mathbf{I}}
\newcommand{\X}{\mathcal{X}}
\newcommand{\gS}{\mathcal{S}}
% \newcommand{\Sm}{\Sigma}
% \newcommand{\La}{\Lambda}



% Figure reference, lower-case.
\def\figref#1{figure~\ref{#1}}
% Figure reference, capital. For start of sentence
\def\Figref#1{Figure~\ref{#1}}
\def\twofigref#1#2{figures \ref{#1} and \ref{#2}}
\def\quadfigref#1#2#3#4{figures \ref{#1}, \ref{#2}, \ref{#3} and \ref{#4}}
% Section reference, lower-case.
\def\secref#1{section~\ref{#1}}
% Section reference, capital.
\def\Secref#1{Section~\ref{#1}}
% Reference to two sections.
\def\twosecrefs#1#2{sections \ref{#1} and \ref{#2}}
% Reference to three sections.
\def\secrefs#1#2#3{sections \ref{#1}, \ref{#2} and \ref{#3}}
% Reference to an equation, lower-case.
\def\eqref#1{equation~\ref{#1}}
% Reference to an equation, upper case
\def\Eqref#1{Equation~\ref{#1}}
\def\TwoEqref#1#2{Equation \ref{#1} and \ref{#2}}
% Reference to an algorithm, lower-case.
\def\algref#1{algorithm~\ref{#1}}
% Reference to an algorithm, upper case.
\def\Algref#1{Algorithm~\ref{#1}}
\def\twoalgref#1#2{algorithms \ref{#1} and \ref{#2}}
\def\Twoalgref#1#2{Algorithms \ref{#1} and \ref{#2}}